% solution-template.tex
\documentclass[11pt,a4paper]{article}
\usepackage[utf8]{inputenc}
\usepackage[T1]{fontenc}
\usepackage[ngerman, english]{babel}
\usepackage{geometry}
\usepackage{amsmath}
\usepackage{fancyhdr}
\usepackage{hyperref}

\geometry{margin=25mm}
\pagestyle{fancy}
\fancyhf{}
\renewcommand{\headrulewidth}{0.4pt}
\lhead{\textbf{Name:} \studentname}
\rhead{\textbf{Matr.-Nr.:} \studentnumber}
\cfoot{\thepage}

% User info - bitte anpassen
\newcommand{\studentname}{Matthias Jan\ss{}en}
\newcommand{\studentnumber}{1871808}
\newcommand{\course}{Grundlagen der Computergrafik}
\newcommand{\institute}{Universit\"at Trier}
\newcommand{\sheetnumber}{05}
\newcommand{\tutor}{Jan Niclas Ruppenthal\\Prof. Dr. Stephan Diehl}

\title{\course}
\author{\studentname\\Matrikelnummer: \studentnumber\\\institute}
\date{\today}

\begin{document}

\maketitle

\begin{center}
    \textbf{Abgabe:} \"Ubungsblatt \sheetnumber \\
    \vspace{4mm}
    \textbf{Tutor:} \tutor
\end{center}

\bigskip

\section*{Aufgaben}
% Hier beginnen deine Loesungen.

\section{Aufgabe 1}
Gegeben (aus der Skizze):
\[
\angle(L_1,\text{Ebene})=20^\circ,\quad
\angle(L_2,\text{Ebene})=80^\circ,\quad
\angle(V_1,\text{Ebene})=75^\circ.
\]
Mit $N \perp$ Ebene:
\[
\angle(L_1,N)=70^\circ,\quad \angle(L_2,N)=10^\circ,
\quad \langle L_1,N\rangle=\cos 70^\circ=\sin 20^\circ,
\quad \langle L_2,N\rangle=\cos 10^\circ=\sin 80^\circ.
\]
Reflexionswinkel für $P_1$:
\[
\angle(R_1,V_1)=75^\circ-20^\circ=55^\circ,
\qquad
\angle(R_2,V_1)=80^\circ-75^\circ=5^\circ.
\]

Kamera $P_1$:
$d_L = 2$ also $f_{at} = \frac{1}{2}$ \\
$I_a^{\mathrm{rot}} = I_a^{\mathrm{gruen}} = I_a^{\mathrm{blau}} = 0.3$.

Rot:
\[
\begin{aligned}
I_{\mathrm{rot}}
&= I_a^{\mathrm{rot}} \cdot k_a^{\mathrm{rot}} \\
&\quad + f_{at}\Bigl(
 I_{L_1}^{\mathrm{rot}} \cdot k_d^{\mathrm{rot}} \cdot \max(0,\langle L_1,N\rangle) \\
&\qquad + I_{L_1}^{\mathrm{rot}} \cdot k_s^{\mathrm{rot}} \cdot \max(0,\langle R_1,V_1\rangle)^{51} \\
&\qquad + I_{L_2}^{\mathrm{rot}} \cdot k_d^{\mathrm{rot}} \cdot \max(0,\langle L_2,N\rangle) \\
&\qquad + I_{L_2}^{\mathrm{rot}} \cdot k_s^{\mathrm{rot}} \cdot \max(0,\langle R_2,V_1\rangle)^{51}
\Bigr) \\
&= 0.3 \cdot 0.33
 + \frac{1}{2}\Bigl(
 0.5 \cdot 0.51 \cdot \sin 20^\circ \\
&\qquad + 0.5 \cdot 0.25 \cdot \cos^{51} 55^\circ
 + 0.75 \cdot 0.51 \cdot \sin 80^\circ \\
&\qquad + 0.75 \cdot 0.25 \cdot \cos^{51} 5^\circ
\Bigr) \\
&= 0.099
 + \frac{1}{2}\Bigl(0.0873 + 0.00000000000012 + 0.3767 + 0.1544\Bigr) \\
&\approx 0.4081.
\end{aligned}
\]

Grün:
\[
\begin{aligned}
I_{\mathrm{gruen}}
&= 0.3 \cdot 0.33
 + \frac{1}{2}\Bigl(
 0.625 \cdot 0.51 \cdot \sin 20^\circ \\
&\qquad + 0.625 \cdot 0.25 \cdot \cos^{51} 55^\circ
 + 0 \cdot 0.51 \cdot \sin 80^\circ
 + 0 \cdot 0.25 \cdot \cos^{51} 5^\circ
\Bigr) \\
&= 0.099
 + \frac{1}{2}\Bigl(0.1090 + 0 + 0 + 0\Bigr) \\
&\approx 0.1535.
\end{aligned}
\]

Blau:
\[
\begin{aligned}
I_{\mathrm{blau}}
&= 0.3 \cdot 0.33
 + \frac{1}{2}\Bigl(
 0 \cdot 0.51 \cdot \sin 20^\circ \\
&\qquad + 0 \cdot 0.25 \cdot \cos^{51} 55^\circ
 + 1 \cdot 0.51 \cdot \sin 80^\circ
 + 1 \cdot 0.25 \cdot \cos^{51} 5^\circ
\Bigr) \\
&= 0.099
 + \frac{1}{2}\Bigl(0 + 0 + 0.5028 + 0.2058\Bigr) \\
&\approx 0.4528.
\end{aligned}
\]

Damit ergeben sich für die drei Farbkanäle
\[
\begin{aligned}
I_{\mathrm{rot}} &\approx 0.408, \\
I_{\mathrm{gruen}} &\approx 0.153, \\
I_{\mathrm{blau}} &\approx 0.453.
\end{aligned}
\]

Sichtbare Farbe der Ebene am Punkt $Q$ aus Sicht von $P_1$:
\[
\bigl(I_{\mathrm{rot}}, I_{\mathrm{gruen}}, I_{\mathrm{blau}}\bigr)
\approx (0.408, 0.154, 0.453).
\]

\section{Aufgabe 2}
Kamera $P_2$:
\[
\angle(V_2,\text{Ebene})=15^\circ,
\quad \angle(R_1,V_2)=20^\circ-15^\circ=5^\circ,
\quad \angle(R_2,V_2)=80^\circ-15^\circ=65^\circ.
\]
$f_{at} = \frac{1}{2}$, $I_a^{\mathrm{rot}} = I_a^{\mathrm{gruen}} = I_a^{\mathrm{blau}} = 0.3$.

Rot:
\[
\begin{aligned}
I_{\mathrm{rot}}
&= 0.3 \cdot 0.33
 + \frac{1}{2}\Bigl(
 0.5 \cdot 0.51 \cdot \sin 20^\circ \\
&\qquad + 0.5 \cdot 0.25 \cdot \cos^{51} 5^\circ
 + 0.75 \cdot 0.51 \cdot \sin 80^\circ \\
&\qquad + 0.75 \cdot 0.25 \cdot \cos^{51} 65^\circ
\Bigr) \\
&= 0.099
 + \frac{1}{2}\Bigl(0.0872 + 0.1029 + 0.3767 + 0\Bigr) \\
&\approx 0.382.
\end{aligned}
\]

Grün:
\[
\begin{aligned}
I_{\mathrm{gruen}}
&= 0.3 \cdot 0.33
 + \frac{1}{2}\Bigl(
 0.625 \cdot 0.51 \cdot \sin 20^\circ \\
&\qquad + 0.625 \cdot 0.25 \cdot \cos^{51} 5^\circ
 + 0 \cdot 0.51 \cdot \sin 80^\circ
 + 0 \cdot 0.25 \cdot \cos^{51} 5^\circ
\Bigr) \\
&= 0.099
 + \frac{1}{2}\Bigl(0.1090 + 0.1286 + 0 + 0\Bigr) \\
&\approx 0.218.
\end{aligned}
\]

Blau:
\[
\begin{aligned}
I_{\mathrm{blau}}
&= 0.3 \cdot 0.33
 + \frac{1}{2}\Bigl(
 0 \cdot 0.51 \cdot \sin 20^\circ \\
&\qquad + 1 \cdot 0.51 \cdot \sin 80^\circ
 + 1 \cdot 0.25 \cdot \cos^{51} 65^\circ
\Bigr) \\
&= 0.099
 + \frac{1}{2}\Bigl(0 + 0.5028 + 0\Bigr) \\
&\approx 0.350.
\end{aligned}
\]

Damit ergeben sich für die drei Farbkanäle
\[
\begin{aligned}
I_{\mathrm{rot}} &\approx 0.382, \\
I_{\mathrm{gruen}} &\approx 0.218, \\
I_{\mathrm{blau}} &\approx 0.350.
\end{aligned}
\]

Sichtbare Farbe der Ebene am Punkt $Q$ aus Sicht von $P_2$:
\[
\bigl(I_{\mathrm{rot}}, I_{\mathrm{gruen}}, I_{\mathrm{blau}}\bigr)
\approx (0.382, 0.218, 0.350).
\]

\end{document}
